% =============================================================================
% Discussion -- draft for P3 (issue #73, closing stress-test Gap B)
% Sources: docs/rfi-questions.md, mrfm/mrfm-v0.1.md ("Limitations of MRFM
% v0.1"), docs/prior-art.md sec 1, GitHub issue #45 (points 3 and 6),
% planning/RESOURCES.md's Track 2 judging criterion (paraphrased below, not
% quoted -- planning/ is gitignored and not present in this worktree). See
% report/report.tex lines 117-125 for the placeholder this replaces and its
% own TODO comment.
% =============================================================================

\section{Discussion}
\label{sec:discussion}

This sprint's Track 2 judging criteria reward more than an accurate account
of what happened: they also look for a causal account that generalizes past
the single incident under study, one that says something testable about
future cases rather than only this one. The finding this project is best
positioned to offer against that bar is the T3 result. Intent-attribution
claims are the one claim type this ledger's own metrics single out as
CoT-dependent -- M1 shows 11.1\% of T3 claims (1 of 9) rest solely on raw
chain-of-thought with no corroborating channel at all, and M2's per-claim-type
breakdown shows T3 claims surviving the loss of chain-of-thought (D1,
opacity) at only 22.2\%, against 89--100\% survival for every other claim
type (Section~\ref{sec:results}). That gap is not an idiosyncrasy of this
incident, this lab, or this model. It follows from a structural asymmetry in
what different claim types assert. An event claim (T1) or a mechanism claim
(T2) asserts something about the world external to the acting system -- a
timestamp, a file touched, an API called -- and such assertions can, in
principle, leave independent traces in channels the agent itself did not
write: host-side logs, cloud-provider audit trails, a second party's own
record. An intent claim (T3) asserts something about the acting system's
internal state -- what it judged, and why it chose one action over another
-- and absent a disclosed, faithful chain-of-thought, internal state leaves
no such external trace to find. It can only be inferred behaviorally, from
what the agent did, and behavioral inference is a strictly lower-fidelity
signal than a contemporaneous, verbalized account of the reasoning itself
would be. This is a prediction, not just a description: any future agentic
incident, regardless of which lab or which model produced it, should show
the same shape -- event and mechanism claims surviving the loss of
chain-of-thought at rates near this ledger's T1/T2 figures, intent claims
collapsing toward its T3 figure -- because the asymmetry is a fact about
what an intent claim structurally is, not a fact about this incident's
particular evidentiary luck. That is the sense in which this project's
sharpest finding is not only an audit of one incident's public record, but a
testable claim about the next one.

This structural asymmetry has a direct corollary once chain-of-thought
itself stops being treated as a fixed feature of future models. The
CoT-monitorability-fragility literature already surveyed in Related Work
(Section~\ref{sec:related-work}) argues that today's externalized,
human-readable reasoning trace is an emergent property of current training
methods, not a guarantee: Korbak et al.'s cross-lab position paper names a
shift toward less monitorable, less verbalized reasoning -- optimizing
directly against a CoT monitor, or moving toward non-linguistic latent
reasoning -- as a live risk, and Baker et al. show a version of this already
happening under direct optimization pressure, where a model taught to evade
a CoT monitor keeps misbehaving while its trace simply stops describing why.
Chen et al.'s hint-based faithfulness tests already put verbalization rates
well below 50\% in some settings even without adversarial pressure on the
monitor. If the trend these sources describe continues, the T3/M1 dependency
finding above does not shrink -- it grows. The one channel intent-attribution
claims currently lean on hardest, imperfectly, is also the channel this
trend degrades fastest, while the alternative -- behavioral inference from
actions alone -- gains no corresponding improvement as a window into internal
state. A world with less monitorable, less verbalized model reasoning is a
world where this ledger's 22.2\% D1-survival rate for T3 claims looks less
like a floor for future incidents' intent claims and more like a ceiling.

The eight RFI questions in \texttt{docs/rfi-questions.md} exist so that this
same channel-fragility finding is not only a description of a problem but
something a regulator -- or a Track 3 team building governance,
verification, or enforcement tooling on top of this sprint's output -- can
act on directly. Each question is anchored to the specific \texttt{claim\_id}s
where the ledger itself, via a \texttt{verifiability=V3} coding or a
\texttt{COLLAPSES} verdict under D1, D2, or D3, already flags the public
record as weakest, rather than to a fresh judgment call made for the RFI
document itself; and each carries a stated ``what answer would settle this''
and ``what answer would NOT settle this'' pair, so a recipient's response can
be checked against a fixed bar instead of read for tone or reassurance
value. That is exactly the shape of question a regulator needs in hand to
actually send, rather than a general request for more transparency, and
exactly the kind of ready-made, evidence-anchored artifact a Track 3 team
could pick up and operationalize without re-deriving it from the raw
125-row ledger -- a concrete, checkable question with a stated pass/fail
bar, which is the other half of what this track's judging criteria reward
alongside the causal claim above.

A month of follow-up would strengthen the method itself, not just this one
application of it. Three additions matter most. First, multi-rater coding
throughout the ledger, rather than the single blind-recode subsample this
project used to compute Cohen's $\kappa$ per degradation condition
(Section~\ref{sec:results}; \texttt{docs/reliability-report.md}) -- a full
second coder on every row would tighten every headline number's confidence
interval, not just confirm that the codebook survived one spot-check.
Second, running the same instrument -- codebook, degradation conditions, and
headline metrics unchanged -- against a second, independent agentic
incident, to test whether the channel-fragility pattern found here, and the
T3-specific prediction above in particular, replicates outside this one
case, or whether this incident's specific mix of disclosed and withheld
evidence was unusual. Third, piloting the RFI questions in
\texttt{docs/rfi-questions.md} against an actual regulatory review process,
or at minimum against feedback from someone who has drafted or received a
real RFI, to check whether the questions as phrased function as intended
once someone outside this project tries to use them. (MRFM v0.1's own
limitations and its own month-of-follow-up list, in
\texttt{mrfm/mrfm-v0.1.md}, are scoped to the fourteen-clause table
specifically and are not repeated here.)

Returning to this project's thesis: the best-powered evidence for it is not
the T3 finding this section leads with structurally, but M5. Across all 125
claims in the ledger, 78.4\% rest on a primary channel the codebook itself
ranks as fragile (C1, C2, C3, or C6), and 47.2\% -- just under half of
everything publicly known about this incident -- rest on such a channel with
no corroborating channel at all. M1's 11.1\% (n=9) is the sharpest, most
viscerally specific illustration of the same underlying pattern, and the
causal argument above is where that pattern's mechanism is easiest to state
precisely, but M5 is the number that should carry this paper's statistical
weight when the claim is about the incident's public record as a whole: not
a fragile statistic about nine claims, but a well-powered one about all of
them.

In sum: this project set out to measure, rather than assume, how much of the
public understanding of one agentic-AI incident rests on evidentiary
channels likely to be absent, unfaithful, or forged next time. The answer is
a large fraction under any of the cuts this ledger supports; the mechanism
behind the sharpest slice of that fraction, intent attribution, is
structural rather than incidental to this incident; and the direction of
travel in how future models reason makes that mechanism more consequential,
not less. MRFM v0.1 (Section~\ref{sec:mrfm}) and the RFI questions above are
this project's answer to what breaks next and what to do before it does;
whether they hold up is exactly the kind of check the follow-up work above
is designed to run.

% -----------------------------------------------------------------------------
% CITE lines -- external works named above beyond those already listed in
% report/drafts/related-work.tex's own CITE block. Full annotations live in
% docs/prior-art.md sec 1; reproduced here for the assembly issue's
% References section.
% -----------------------------------------------------------------------------
% CITE: Tomek Korbak, Mikita Balesni, Elizabeth Barnes, et al. (cross-lab position paper), "Chain of Thought Monitorability: A New and Fragile Opportunity for AI Safety," July 2025. https://arxiv.org/abs/2507.11473
% CITE: Bowen Baker, Joost Huizinga, Leo Gao, Zehao Dou, Melody Y. Guan, Aleksander Madry, Wojciech Zaremba, Jakub Pachocki, David Farhi (OpenAI), "Monitoring Reasoning Models for Misbehavior and the Risks of Promoting Obfuscation," March 2025. https://arxiv.org/abs/2503.11926
% CITE: Yanda Chen, Joe Benton, Ansh Radhakrishnan, Jonathan Uesato, Carson Denison, John Schulman, Arushi Somani, et al. (Anthropic), "Reasoning Models Don't Always Say What They Think," May 2025. https://arxiv.org/abs/2505.05410
